\begin{definition}
Fix a finite loopless multigraph \(G\), a color bound \(k\), and a
distinguished edge \(e_0\).  A \emph{\(T\)-change counterexample} consists
of a partial coloring \(\varphi\), a tree-sequence \(T\), two distinct
colors \(a,b\), three outside vertices \(a_1,a_2,a_3\), three tree
vertices \(b_1,b_2,b_3\), three incidence-graph paths \(P_1,P_2,P_3\),
and three exact edges \(f_1,f_2,f_3\), subject to all of the following.
The tuple \((G,e_0,\varphi)\) is a \(k\)-triple; \(T\) is a closed
Tashkinov tree; \(a\) is missing on \(T\), while \(b\) is not missing on
\(T\).  For each \(i\), \(P_i\) is a
\(\noderef{TreeExitPath}\) in the \((a,b)\)-incidence graph from \(a_i\)
to \(b_i\).  The three paths are pairwise support-disjoint, the three
outside ends are pairwise distinct, and the three tree ends are pairwise
distinct.  Finally, \(f_i\) is an exact \(b\)-colored boundary edge of
\(V(T)\), is incident with \(b_i\), and its edge-node belongs to the
support of \(P_i\).
\end{definition}
